\begin{table}[htbp]
\centering
\caption{E2 --- Solver benchmark on three representative instances (5 replications each). Time limit: 3\,600\,s.}
\label{tab:e2_benchmark}
\scriptsize
\begin{tabular}{llcrrrr}
\toprule
\textbf{Instance} & \textbf{Solver} & \textbf{Status}
  & \textbf{Time (s) mean} & \textbf{Std (s)} & \textbf{Obj.\ mean} & \textbf{Opt.\ rate (\%)} \\
\midrule
  \multirow{4}{*}{S1 (large)} & Gurobi & Optimal & 128.1 & 3.3 & 333.0 & 100 \\
   & HiGHS & Optimal & 1040.3 & 1.7 & 333.0 & 100 \\
   & CBC & Feas. (0\% opt) & 3593.3 & 3.9 & 342.0 & 0 \\
   & GLPK & Feas. (0\% opt) & 6599.8 & 10.8 & 411.0 & 0 \\
\midrule
  \multirow{4}{*}{S3 (medium)} & Gurobi & Optimal & 64.5 & 1.8 & 654.0 & 100 \\
   & HiGHS & Optimal & 524.7 & 1.8 & 654.0 & 100 \\
   & CBC & Feas. (0\% opt) & 3677.1 & 7.6 & 997.0 & 0 \\
   & GLPK & Feas. (0\% opt) & 5398.0 & 5.8 & 280.0 & 0 \\
\midrule
  \multirow{4}{*}{S9 (small)} & Gurobi & Optimal & 12.5 & 6.3 & 43.0 & 100 \\
   & HiGHS & Optimal & 36.7 & 2.0 & 43.0 & 100 \\
   & CBC & Optimal & 127.1 & 7.0 & 43.0 & 100 \\
   & GLPK & Feas. (0\% opt) & 1819.0 & 1.4 & 43.0 & 0 \\
\botrule
\end{tabular}
\begin{tablenotes}
\small
\item Opt.\ rate = percentage of replications reaching proven optimality within the time limit.
\item All solvers use the same model configuration as E1 (huecos\_grupo=True, preferencias=True).
\end{tablenotes}
\end{table}
